{\scriptsize
        \begin{tabularx}{0.55\textwidth}{|p{0.18\columnwidth}|X|}
            \hline 
            \textbf{Notions et contenus} & \textbf{Capacités exigibles} \\
            \hline 
            \multicolumn{2}{|l|}{\textbf{ 6.1. Phénomènes de propagation non dispersifs : équation de d'Alembert }} \\
            \hline
            \multicolumn{2}{|l|}{\textbf{
            6.1.1. Ondes mécaniques unidimensionnelles dans les solides déformables }}\\
            \hline
            Ondes transversales sur une corde vibrante.  & 
            Établir l'équation d'onde décrivant les ondes transversales sur une corde vibrante infiniment souple dans l'approximation des petits mouvements transverses. \\
            \hline 
            Domaine d'élasticité d'un solide : module d'Young, loi de Hooke.   & 
            Exploiter le modèle de la chaîne d'atomes élastiquement liés pour relier le module d'Young d'un solide élastique à ses caractéristiques microscopiques. \\
        \hline 
        Ondes mécaniques longitudinales dans une tige solide dans l'approximation des milieux continus. & 
        Établir l'équation d'onde décrivant les ondes mécaniques longitudinales dans une tige solide. \\ 
        \hline
        \end{tabularx}
        \begin{tabularx}{0.43\textwidth}{|p{0.15\columnwidth}|X|}
            \hline
            Équation de d'Alembert ; célérité.  & 
            Identifier l'équation de d'Alembert. Relier qualitativement la célérité d'ondes mécaniques, la raideur et l'inertie du milieu support. \\
            \hline 
            Ondes progressives, ondes progressives harmoniques ; ondes stationnaires.   & 
            Différencier une onde stationnaire d'une onde progressive.
Utiliser qualitativement l'analyse de Fourier pour décrire une onde non harmonique. \\
        \hline 
        Modes propres d'une corde vibrante fixée à ses deux extrémités. Résonances d'une corde de Melde. & 
        Décrire les modes propres d'une corde vibrante fixée à ses deux extrémités. Interpréter quantitativement les résonances observées avec la corde de Melde en négligeant l'amortissement. \\ 
        \hline
        \end{tabularx}
\vspace{0.5cm}
}
